\section{Theoretical Analysis}
\label{sec:theoretical-analysis}
Our theoretical analysis (cf. \textbf{\Cref{appendix:theoretical-analysis}} for a detailed proof) of the memory capacity demonstrates that, given a set of queries exhibiting some degree of semantic similarity, our memorization rules in \Cref{eq:update-function} enable edge embeddings to memorize these queries and store sufficient information, thereby enabling effective memory replay. Formally, when the embedding dimension $d$ is sufficiently large, given the threshold $\lambda$ in Section \ref{sec::memory-replay}, for any given edge, when the maximum angle $\theta$ between all query embedding pairs within a query set satisfies \Cref{eq:lamba-theta}, the edge can effectively memorize the semantic information within that query set. 
\begin{equation}
    \theta \leq \lim_{d \to \infty} \left[ 2\arcsin\left(\sqrt{\frac{1}{2}}\sin(\arccos(\lambda))\right) \right]
\label{eq:lamba-theta}
\end{equation}



\begin{remark}
    To satisfy the storage requirements\footnote{While this provides a solution under theoretical conditions for parameter configuration, practical applications require consideration of multiple factors when setting parameters. Additional details are in \Cref{appendix:parameter-configuration}.}, the theoretical upper bound for $\lambda$ is determined to be 0.775. This determination aligns with existing research \cite{Base1}, which establishes a cosine similarity threshold of 0.6 as the criterion for assessing semantic similarity.
\end{remark}

\begin{remark}
Our assumption that the embedding dimension $d$ is infinitely large is based on the fact that linear embedding models do have very high dimensions. For the function $\sqrt{\frac{d+1}{2d}}$, when $d$ is greater than around 100 (currently, the embeddings of an LLM often exceed 100 dimensions), it can actually be approximated to $\frac{1}{2}$. This approximation thus holds significant practical utility. See our empirical experiments in Section \ref{ms}.
\end{remark}